\section*{Capítulo \ref{ch:b1:ecosystem-organization} --- Ecossistemas e estrutura trófica}

\begin{solution}{exo:b1:ecosystem-organization:1}
O \omterm{def:b1:ecosystem-organization:ecosystem}{ecossistema} da lagoa é a \omterm{def:b1:ecosystem-organization:ecosystem}{comunidade} --- todas as suas \omterm{def:b1:populations:population}{populações}:
algas, elódeas, caramujos, larvas de insetos, rãs, peixes, bactérias
--- junto com o \omterm{def:b1:ecosystem-organization:ecosystem}{biótopo}: a água, sua temperatura, a luz, os gases
dissolvidos e os minerais, o lodo.
\end{solution}

\begin{solution}{exo:b1:ecosystem-organization:2}
\omterm{def:b1:ecosystem-organization:niche}{Hábitat}: o lugar onde a espécie vive (uma floresta de abetos). \omterm{def:b1:ecosystem-organization:niche}{Nicho}:
seu modo de viver ali --- o que come, quando, em que parte da árvore,
o que tolera. \omterm{def:b1:ecosystem-organization:niche}{Nicho} fundamental: a faixa que ela poderia ocupar
sozinha; realizado: o que os competidores e os predadores lhe deixam.
\end{solution}

\begin{solution}{exo:b1:ecosystem-organization:3}
$20\,800/1.7\times 10^6 = \qty{1.2}{\%}$; os herbívoros tomaram
$3400/8800 = \qty{39}{\%}$ da produção líquida.
\end{solution}

\begin{solution}{exo:b1:ecosystem-organization:4}
\omterm{def:b1:ecosystem-organization:trophic}{Produtores}, \omterm{def:b1:ecosystem-organization:trophic}{consumidores} (primários, secundários, terciários),
\omterm{def:b1:ecosystem-organization:trophic}{decompositores}, detritívoros; os \omterm{def:b1:ecosystem-organization:trophic}{decompositores}, que devolvem os
minerais da matéria morta aos \omterm{def:b1:ecosystem-organization:trophic}{produtores}.
\end{solution}

\begin{solution}{exo:b1:ecosystem-organization:5}
PPL $= 800\times 17 = \qty{13.6}{MJ/m^2}$; comido $\qty{2.04}{MJ}$;
assimilado $\qty{1.02}{MJ}$; coelho $\qty{51}{kJ/m^2}$, isto é,
\qty{510}{MJ} por hectare (\qty{30}{kg} de coelho a \qty{17}{kJ/g}).
Eficiência $51/13\,600 = \qty{0.4}{\%}$.
\end{solution}

\begin{solution}{exo:b1:ecosystem-organization:6}
A biomassa é um estoque: um fitoplâncton comido tão depressa quanto
cresce não precisa acumular-se, e por isso o estoque presente de
\omterm{def:b1:ecosystem-organization:trophic}{produtores} pode ser menor do que o dos animais de vida mais longa que
dele se alimentam. A energia é um fluxo: cada nível recebe do de baixo
tudo o que jamais terá e tem de perder parte disso como calor, de modo
que o fluxo só pode encolher para cima.
\end{solution}

\begin{solution}{exo:b1:ecosystem-organization:7}
$p = 0.5, 0.3, 0.15, 0.05$: $H = -(0.5\ln 0.5 + 0.3\ln 0.3 + 0.15\ln
0.15 + 0.05\ln 0.05) = 0.347 + 0.361 + 0.285 + 0.150 = 1.14$;
equitabilidade $1.14/\ln 4 = 0.82$. Quatro com 25 cada: $H = \ln 4 = 1.39$,
equitabilidade 1.
\end{solution}

\begin{solution}{exo:b1:ecosystem-organization:8}
Líquida $= \qty{3}{mg}$ de $\mathrm{O_2}$ por litro por dia;
respiração \qty{1}{mg}; bruta $3 + 1 = \qty{4}{mg}$. Em carbono: bruta
\qty{1.5}{mg}, líquida \qty{1.1}{mg} de C por litro por dia.
\end{solution}

\begin{solution}{exo:b1:ecosystem-organization:9}
Cada elo guarda um décimo: dez elos deixariam um bilionésimo da
produção primária, pouco demais para alimentar um único animal em
qualquer área. Onde a produção é baixa (deserto) ou a área é pequena
(ilha), o décimo de um décimo se esgota depois de menos elos; a
produção alta de uma floresta sustenta um ou dois a mais.
\end{solution}

\begin{solution}{exo:b1:ecosystem-organization:10}
Grão: \qty{1}{t} exige um sexto de hectare. Carne bovina: \qty{1}{t}
de energia em carne exige \qty{10}{t} de grão, \num{1.7} hectares ---
dez vezes a terra. Comer no segundo nível custa o equivalente a um
nível de energia; um planeta alimenta dez vezes mais gente com grão do
que com carne alimentada a grão, e o \omterm{def:b1:ecosystem-organization:trophic}{nível trófico} humano, na média
das dietas, é um termo importante da \omterm{thm:b1:populations:logistic}{capacidade de suporte}.
\end{solution}

\begin{solution}{exo:b1:ecosystem-organization:11}
Por metro quadrado o oceano é um deserto, mas cobre dois terços do
planeta, e por isso uma taxa baixa vezes uma área enorme dá metade do
total. Sua produção é limitada pelos nutrientes, que afundam para fora
da camada iluminada; onde as correntes os trazem de volta à superfície
(ressurgências costeiras, plataformas), a produção é dez vezes maior,
e é aí, em poucos por cento do mar, que se concentram os peixes e as
pescarias.
\end{solution}

\begin{solution}{exo:b1:ecosystem-organization:12}
A energia entra como luz, atravessa alguns níveis de \omterm{def:b1:organism-environment:organism}{organismos} e sai
inteiramente como calor: nesse sentido o \omterm{def:b1:ecosystem-organization:ecosystem}{ecossistema} de fato
transforma luz solar em calor, e os \omterm{def:b1:ecosystem-organization:trophic}{decompositores}, que respiram tudo
o que os \omterm{def:b1:ecosystem-organization:trophic}{consumidores} não respiraram, são o destino da maior parte
dela. Mas a matéria não é consumida: carbono, nitrogênio e fósforo
circulam, devolvidos aos \omterm{def:b1:ecosystem-organization:trophic}{produtores} pelos mesmos \omterm{def:b1:ecosystem-organization:trophic}{decompositores}, e por
isso o sistema persiste; e os \omterm{def:b1:organism-environment:organism}{organismos} não são um subproduto, mas a
estrutura que organiza o fluxo --- os \omterm{def:b1:ecosystem-organization:niche}{nichos}, as teias e as pirâmides
são o que a energia constrói ao passar. A frase acerta a termodinâmica
e perde a biologia.
\end{solution}

\begin{solution}{pb:b1:ecosystem-organization:1}
\textbf{1.} $20\,800 - 12\,000 = \qty{8800}{kcal}$ por metro quadrado
por ano.
\textbf{2.} \qty{1.2}{\%}. Perdas: luz não absorvida (metade do
espectro, reflexão, água), \omterm{def:b1:flowering-plant-organization:organs}{folhas} saturadas em pleno sol ou
sombreadas, fotorrespiração, e a respiração das próprias plantas (já
contada entre o bruto e o líquido).
\textbf{3.} $12\,000/20\,800 = \qty{58}{\%}$.
\textbf{4.} Comido \qty{3400}{}; diretamente para os \omterm{def:b1:ecosystem-organization:trophic}{decompositores},
$8800 - 3400 = \qty{5400}{kcal}$.
\textbf{5.} $8800/10 = \qty{880}{g/m^2}$: como uma floresta temperada
ou um campo rico.
\textbf{6.} $1.7\times 10^6 - 20\,800 = \qty{1.68e6}{kcal}$, \qty{98.8}{\%}
do total: refletida, transmitida, ou absorvida pela água e pela rocha
e reirradiada como calor.
\textbf{7.} Produção $1500 - 1100 = \qty{400}{kcal}$; eficiência de
assimilação $1500/3400 = \qty{44}{\%}$.
\textbf{8.} $400/1500 = \qty{27}{\%}$.
\textbf{9.} $400/8800 = \qty{4.5}{\%}$ (\qty{16}{\%} se calculada sobre
o que foi comido, \qty{3400}{}).
\textbf{10.} Produção dos carnívoros $380 - 300 = \qty{80}{kcal}$;
eficiência $80/400 = \qty{20}{\%}$ (eles comeram \qty{380}{} dos
\qty{400}{} produzidos).
\textbf{11.} $6/80 = \qty{7.5}{\%}$.
\textbf{12.} \qty{6}{kcal}; $6/1.7\times 10^6 = \num{3.5e-6}$, alguns
milionésimos.
\textbf{13.} A carne é digerível e próxima em composição do que a
come; a matéria vegetal é fibrosa e em boa parte indigerível, de modo
que um herbívoro perde nas fezes uma parte maior do que ingere.
\textbf{14.} Produção líquida não comida \qty{5400}{}; fezes dos
herbívoros $3400 - 1500 = \qty{1900}{}$; produção dos herbívoros não
comida $400 - 380 = \qty{20}{}$; fezes dos carnívoros (a assimilação
não é dada; tome comido menos respirado menos produção $= 0$, tudo
assimilado) e produção dos carnívoros não comida $80 - 20 = \qty{60}{}$;
produção dos carnívoros de topo \qty{6}{}: total \qty{7386}{kcal}.
\textbf{15.} $12\,000 + 1100 + 300 + 14 + 7386 = \qty{20800}{kcal}$.
\textbf{16.} Igual à produção bruta: a nascente está em estado
estacionário e não armazena biomassa de um ano para o outro.
\textbf{17.} $7386/20\,800 = \qty{36}{\%}$ (a maior parte do resto é a
respiração dos próprios \omterm{def:b1:ecosystem-organization:trophic}{produtores}).
\textbf{18.} \omterm{def:b1:ecosystem-organization:trophic}{Produtores} \qty{8800}{}, herbívoros \qty{400}{},
carnívoros \qty{80}{}, carnívoros de topo \qty{6}{}; eficiências
\qty{4.5}{\%}, \qty{20}{\%}, \qty{7.5}{\%}.
\textbf{19.} \qty{8000}{}, \qty{400}{}, \qty{100}{}, \qty{15}{kcal/m^2}:
direita, cada nível um décimo ou menos do de baixo.
\textbf{20.} \omterm{def:b1:ecosystem-organization:trophic}{Produtores} $8000/8800 = 0.9$ ano; herbívoros $400/400 =
1$ ano; carnívoros $100/80 = 1.25$ ano; carnívoros de topo $15/6 =
2.5$ anos. As plantas se renovam mais depressa: \omterm{def:b1:body-plans-tissues:tissue}{tecido} de vida curta,
pastado continuamente; os carnívoros de topo são animais de vida longa
cujo estoque se renova devagar.
\textbf{21.} Uma garça de \qty{2}{kg} tem cerca de \qty{500}{g} secos,
\qty{5000}{kcal}, e precisa aproximadamente disso em produção a cada
ano para repor a si mesma e à sua prole: $5000/6 \approx \qty{800}{m^2}$
de nascente --- na prática muito mais, pois ela também respira: a
\qty{100}{kcal} por dia, \qty{36500}{kcal} por ano de ingestão, isto
é, \qty{6000}{m^2} do nível abaixo dela (produção de carnívoros
\qty{80}{}), o que corresponde ao território de uma garça.
\textbf{22.} A produção líquida excedente não é comida e vai para os
\omterm{def:b1:ecosystem-organization:trophic}{decompositores}, cuja respiração duplica; à noite, sem fotossíntese, a
demanda de oxigênio deles pode esvaziar a água de oxigênio e matar os
peixes --- a eutrofização.
\textbf{23.} A produção dos herbívoros cai para \qty{200}{}: os
carnívoros recebem metade do alimento, sua produção cai rumo a
\qty{40}{}, a dos carnívoros de topo rumo a 3, e há menos garças; a
erva-marinha, menos pastada, adensa-se e uma parte maior dela vai sem
ser comida para os \omterm{def:b1:ecosystem-organization:trophic}{decompositores}.
\textbf{24.} A energia dos carnívoros de topo é ínfima, mas seu efeito
recai sobre o número de carnívoros, que controla os herbívoros, que
controlam o capim: retirá-los deixa os carnívoros subirem, os
herbívoros caírem e o capim crescer --- uma cascata através da teia
(\cref{ch:b1:species-interactions}); retirar algumas quilocalorias de
herbívoros não muda nada além de algumas quilocalorias.
\textbf{25.} \omterm{def:b1:ecosystem-organization:trophic}{Produtores} a herbívoros \qty{4.5}{\%} (da PPL;
\qty{16}{\%} do que foi comido), herbívoros a carnívoros \qty{20}{\%},
carnívoros a carnívoros de topo \qty{7.5}{\%}; luz solar a produção
bruta \qty{1.2}{\%}; alguns milionésimos da luz solar terminam nos
carnívoros de topo.
\end{solution}
